Index calculus is the method that, on paper, should exploit the most
structure in an elliptic-curve discrete-logarithm instance, and in practice
has never beaten Pollard rho on a prime-order curve of cryptographic shape.
Between July and mid-October 2026 we ran a sustained engineering and
measurement programme on this question in one open repository: about
6{,}900 commits and 450 dated study directories, the index-calculus part of
which began in the first week of September. This paper is the record of that
programme. We describe the routes we pursued---Frobenius-invariant factor
bases on Koblitz curves and their compact-orbit form; SAT and Gr\"obner-basis
decomposition oracles (matrix F4, F5, an inherited F4 and an F6-IC variant);
root-finding decomposition oracles; Kummer-tower and symmetry-reduced prime
and extension-field systems; isogeny and transfer routes; and a
twelve-workstream replication of the 2025--26 literature---the
infrastructure that made the results comparable (\emph{ecbench}, a native
harness that charges every ECDLP method in one counted unit, group-addition
equivalents per $\sqrt{r}$, with sealed records, exact replays, isolation
levels and paired cluster-bootstrap intervals), and what improved.

The largest single change came from a suggestion in a manuscript we
audited: the roots of Semaev's third summation polynomial in its free
variable are $x(R+F)$ and $x(R-F)$, so a two-summand decomposition needs no
polynomial system at all. We proved this (also in Lean), generalised it to
meet-in-the-middle pair-sum tables for three and four summands, to a CUDA
kernel, to Kummer-tower bases and to dense root indices, and built it into
the harness. Per decomposition it is $383$--$1{,}535\times$ faster than our
matrix F4 on Koblitz curves and $3{,}021$--$5{,}802\times$ faster than
Buchberger on the five NIST prime curves; end to end it is
$50$--$105\times$ faster than msolve on small prime curves. Inside ecbench
it moved the best complete-count Koblitz index-calculus arm at $n=13$ from
$S=38.8$ to $S=23.0$, below an XOR-priced F4 at $S=421$ but still
$1.55\times$ the strong signed-Frobenius rho. FGHR's $D_4$ symmetry cut the
$k=4$ Gaudry solve over $\FF_{p^4}$ by $47.5$--$48.4\times$, moving a
modelled crossover with rho from $2^{150}$ to $2^{135}$; our $\FF_p$ F4
reached parity with msolve on square systems; and a compact-orbit Koblitz
index calculus with retained precomputation solves one unseen public target
at $n=61$ $13.9\times$ faster online than rho (95\% interval $9.8$--$19.1$,
64 targets, pinned host) while costing $83\times$ rho when its set-up is
charged. We report the negative results with the same weight: no measured
index-calculus arm beats matched rho on the cold, one-target,
set-up-inclusive boundary; a group-homomorphic transfer cannot change
decomposition counts (we state the conservation lemma); and several
literature claims did not survive implementation. Every number cites a
committed artefact; nothing here concerns a deployed curve or key recovery.
